% Provenance (never printed): NSI terminale "Structures de donnees: listes, piles, files; interface et
% implementation; structures lineaires". See sources/book2/BRIEFS.md.
\chapter{Linked Lists, Stacks and Queues}\label{ch:g12:lists-stacks-queues}

The undo button of an editor cancels the last change first, then the one before it: the newest change leaves
first. A printer serves its waiting documents in the order they came: the oldest leaves first. The two are the same
idea, a line of things waiting, taken from opposite ends; and each has a name, an interface, and more than one way to
build it.

This chapter separates what a structure does from how it is built. It builds the \omterm{def:g12:lists-stacks-queues:linked-list}{linked list}, a chain of small \omterm{def:g12:lists-stacks-queues:linked-list}{cells},
and from it the \omterm{def:g12:lists-stacks-queues:stack}{stack} and the \omterm{def:g12:lists-stacks-queues:queue}{queue}, each in two ways; it counts what each operation costs; and it uses them to check
brackets, to compute, and to simulate a day at a post office counter.

\begin{recall}
A class groups data with its methods, and an invariant is kept by the methods (chapter
\ref{ch:g12:classes-and-objects}). A module offers an interface and hides its \omterm{def:g11:modules-and-interfaces:interface}{implementation} (chapter
\ref{ch:g11:modules-and-interfaces}). The cost of an algorithm (chapter \ref{ch:g11:cost-of-an-algorithm}). A
dictionary finds a value from its key (Book 1, grade 8).
\end{recall}

\section{A structure is defined by its operations}\label{sec:g12:lists-stacks-queues:adt}

\begin{definition}[Abstract data type]\label{def:g12:lists-stacks-queues:adt}
An \emph{abstract data type}\index{abstract data type} is a kind of data defined only by its operations and what
each one does: an interface with its specifications. How the data is stored, its \omterm{def:g11:modules-and-interfaces:interface}{implementation}, is left out; one
abstract data type can have several \omterm{def:g11:modules-and-interfaces:interface}{implementations}, and a program that uses only the interface works with any of
them.
\end{definition}

This is the idea of the module interface, applied to data. The \omterm{def:g12:lists-stacks-queues:stack}{stack} and the \omterm{def:g12:lists-stacks-queues:queue}{queue} below are \omterm{def:g12:lists-stacks-queues:adt}{abstract data types}; the
\omterm{def:g12:lists-stacks-queues:linked-list}{linked list} is one way to implement them.

\section{The linked list}\label{sec:g12:lists-stacks-queues:linked}

\begin{definition}[Linked list, cell]\label{def:g12:lists-stacks-queues:linked-list}
A \emph{linked list}\index{linked list} is a sequence stored as a chain of \emph{cells}\index{cell (linked list)}:
each cell holds one value and a link to the next cell, and the last cell's link is \code{None}. The list is known by
its first cell, its head.
\end{definition}

\omcode{code/grade-12/03-lists-stacks-queues/listings/cell.py}{1}{23}{\omterm{def:g12:lists-stacks-queues:linked-list}{Cells}, and two conversions to and from Python
lists.}\label{lst:g12:lists-stacks-queues:cell}

\begin{omfigure}
\begin{tikzpicture}[font=\footnotesize, val/.style={draw, minimum width=7mm, minimum height=6mm, fill=omDef!10},
    lnk/.style={draw, minimum width=5mm, minimum height=6mm}]
  \node at (-1.3,0) {\code{head}};
  \foreach \v [count=\i] in {3, 1, 4, 1} {
    \node[val] (v\i) at ({2.0*(\i-1)},0) {$\v$};
    \node[lnk, right=0pt of v\i] (l\i) {};
  }
  \draw[->, thick] (-0.9,0) -- (v1.west);
  \foreach \i/\j in {1/2, 2/3, 3/4} \draw[->, thick] (l\i.center) -- (v\j.west);
  \fill (l1.center) circle (1.2pt); \fill (l2.center) circle (1.2pt); \fill (l3.center) circle (1.2pt);
  \draw (l4.south west) -- (l4.north east);
  \node[right] at (7.3,0) {\code{None}};
  \node[font=\scriptsize, black!60] at (3.2,-0.8) {each cell: a value, and a link to the next cell};
\end{tikzpicture}

{\small The \omterm{def:g12:lists-stacks-queues:linked-list}{linked list} $3, 1, 4, 1$. Walking it means following the links from the head; the crossed box is the link
\code{None} of the last cell. Inserting at the head is one new cell whose link is the old head.}
\end{omfigure}

\omcode{code/grade-12/03-lists-stacks-queues/listings/linked_demo.py}{1}{15}{Walking, inserting at the head, and
removing a cell by skipping it.}\label{lst:g12:lists-stacks-queues:demo}

\omoutput{code/grade-12/03-lists-stacks-queues/outputs/linked_demo.txt}{Inserting and removing changed one link each;
no value moved.}\label{out:g12:lists-stacks-queues:demo}

A \omterm{def:g12:lists-stacks-queues:linked-list}{linked list} and a Python list cost differently. Reading the value at index $i$ of a Python list takes the same
short time for every $i$; in a \omterm{def:g12:lists-stacks-queues:linked-list}{linked list} it means walking $i$ links. Inserting at the front of a Python list shifts
every element by one place; in a \omterm{def:g12:lists-stacks-queues:linked-list}{linked list} it is one new cell, whatever the length.

\section{The stack}\label{sec:g12:lists-stacks-queues:stack}

\begin{definition}[Stack, push, pop, LIFO]\label{def:g12:lists-stacks-queues:stack}
A \emph{stack}\index{stack} is an \omterm{def:g12:lists-stacks-queues:adt}{abstract data type} holding a sequence of elements, with four operations: create an
empty stack; test whether it is empty; \emph{push}\index{push} an element on top; \emph{pop}\index{pop} the top
element, removing and returning it, which is refused on an empty stack. A fifth, \code{top}, returns the top without
removing it. The last element pushed is the first popped: the stack is \emph{LIFO}\index{LIFO}, last in, first out.
\end{definition}

\begin{omfigure}
\begin{tikzpicture}[font=\footnotesize, el/.style={draw, minimum width=12mm, minimum height=5mm, fill=omDef!10}]
  \foreach \x/\n/\lab in {0/2/{before}, 3.4/3/{after \code{push("c")}}, 6.8/2/{after \code{pop()}}} {
    \draw[thick] (\x-0.75,1.9) -- (\x-0.75,-0.05) -- (\x+0.75,-0.05) -- (\x+0.75,1.9);
    \node[el] at (\x,0.25) {\code{"a"}}; \node[el] at (\x,0.75) {\code{"b"}};
    \ifnum\n=3 \node[el, fill=omThm!15] at (\x,1.25) {\code{"c"}}; \fi
    \node[below] at (\x,-0.1) {\lab};
  }
  \draw[->, omThm, thick] (1.0,1.6) to[bend left] node[above] {push} (2.55,1.35);
  \draw[->, omThm, thick] (4.25,1.35) to[bend left] node[above] {pop returns \code{"c"}} (5.95,1.6);
\end{tikzpicture}

{\small A \omterm{def:g12:lists-stacks-queues:stack}{stack} before and after one \omterm{def:g12:lists-stacks-queues:stack}{push} and one \omterm{def:g12:lists-stacks-queues:stack}{pop}. Only the top is reachable; the \omterm{def:g12:lists-stacks-queues:stack}{pop} gives back the element of
the last \omterm{def:g12:lists-stacks-queues:stack}{push}.}
\end{omfigure}

Two \omterm{def:g11:modules-and-interfaces:interface}{implementations} share one interface. In the first, the top is the end of a Python list; in the second, the top
is the head of a \omterm{def:g12:lists-stacks-queues:linked-list}{linked list}. Both \omterm{def:g12:lists-stacks-queues:stack}{push} and \omterm{def:g12:lists-stacks-queues:stack}{pop} in a time that does not depend on the number of elements.

\omcode{code/grade-12/03-lists-stacks-queues/listings/stack_list.py}{1}{22}{A \omterm{def:g12:lists-stacks-queues:stack}{stack} in a Python
list.}\label{lst:g12:lists-stacks-queues:stack-list}

\omcode{code/grade-12/03-lists-stacks-queues/listings/stack_linked.py}{1}{25}{The same interface, with linked
\omterm{def:g12:lists-stacks-queues:linked-list}{cells}.}\label{lst:g12:lists-stacks-queues:stack-linked}

Since both promise the same behaviour, one \omterm{def:g10:specifications-and-tests:test-suite}{test suite} checks both: the tests use only the interface, and the class
under test is chosen by a loop.

\omcode{code/grade-12/03-lists-stacks-queues/listings/stacks_tests.py}{1}{31}{One suite of tests for the \omterm{def:g12:lists-stacks-queues:stack}{stack}
interface.}\label{lst:g12:lists-stacks-queues:stacks-tests}

\omcode{code/grade-12/03-lists-stacks-queues/listings/stacks_tests.py}{34}{43}{The suite run on each \omterm{def:g11:modules-and-interfaces:interface}{implementation} in
turn.}\label{lst:g12:lists-stacks-queues:stacks-tests-run}

\omoutput{code/grade-12/03-lists-stacks-queues/outputs/stacks_tests.txt}{Both \omterm{def:g11:modules-and-interfaces:interface}{implementations} pass the four
tests.}\label{out:g12:lists-stacks-queues:stacks-tests}

\section{The queue}\label{sec:g12:lists-stacks-queues:queue}

\begin{definition}[Queue, enqueue, dequeue, FIFO]\label{def:g12:lists-stacks-queues:queue}
A \emph{queue}\index{queue} is an \omterm{def:g12:lists-stacks-queues:adt}{abstract data type} holding a sequence, with the operations: create an empty queue;
test whether it is empty; \emph{enqueue}\index{enqueue} an element at the back; \emph{dequeue}\index{dequeue} the
element at the front, removing and returning it, refused on an empty queue. The first element enqueued is the first
dequeued: the queue is \emph{FIFO}\index{FIFO}, first in, first out.
\end{definition}

A \omterm{def:g12:lists-stacks-queues:linked-list}{linked list} makes a \omterm{def:g12:lists-stacks-queues:queue}{queue} if the program keeps a second pointer, to the last cell: elements leave at the head and
arrive after the tail. Both operations then change one or two links.

\begin{omfigure}
\begin{tikzpicture}[font=\footnotesize, val/.style={draw, minimum width=7mm, minimum height=6mm, fill=omDef!10},
    lnk/.style={draw, minimum width=5mm, minimum height=6mm}]
  \foreach \v [count=\i] in {5, 8, 2} {
    \node[val] (v\i) at ({2.0*(\i-1)},0) {$\v$};
    \node[lnk, right=0pt of v\i] (l\i) {};
  }
  \node[val, dashed, fill=omThm!10] (v4) at (6,0) {$7$}; \node[lnk, dashed, right=0pt of v4] (l4) {};
  \foreach \i/\j in {1/2, 2/3} {\fill (l\i.center) circle (1.2pt); \draw[->, thick] (l\i.center) -- (v\j.west);}
  \fill[omThm] (l3.center) circle (1.2pt); \draw[->, thick, dashed, omThm] (l3.center) -- (v4.west);
  \draw (l4.south west) -- (l4.north east);
  \node (h) at (0,1.3) {\code{_head}}; \draw[->, thick] (h) -- (v1.north);
  \node (t) at (4,1.3) {\code{_tail}}; \draw[->, thick] (t) -- (v3.north);
  \draw[->, thick, dashed, omThm] (t.east) to[bend left=20] (v4.north);
  \node[font=\scriptsize, omThm] at (6.3,-0.8) {enqueue 7};
  \node[font=\scriptsize] at (0,-0.8) {dequeue here};
\end{tikzpicture}

{\small A linked \omterm{def:g12:lists-stacks-queues:queue}{queue} holding $5, 8, 2$. Enqueueing $7$ links a new cell after the tail and moves \code{_tail} to it
(dashed); dequeueing moves \code{_head} to the next cell.}
\end{omfigure}

\omcode{code/grade-12/03-lists-stacks-queues/listings/queue_linked.py}{1}{31}{A linked \omterm{def:g12:lists-stacks-queues:queue}{queue}, which also counts its
moves.}\label{lst:g12:lists-stacks-queues:queue-linked}

The second \omterm{def:g11:modules-and-interfaces:interface}{implementation} surprises: a \omterm{def:g12:lists-stacks-queues:queue}{queue} made of two \omterm{def:g12:lists-stacks-queues:stack}{stacks}. Arrivals are pushed on an \omterm{def:g12:lists-stacks-queues:stack}{in-stack}. To \omterm{def:g12:lists-stacks-queues:queue}{dequeue}, the
program \omterm{def:g12:lists-stacks-queues:stack}{pops} from an \omterm{def:g12:lists-stacks-queues:stack}{out-stack}; when the \omterm{def:g12:lists-stacks-queues:stack}{out-stack} is empty, it first \omterm{def:g12:lists-stacks-queues:stack}{pops} every element of the \omterm{def:g12:lists-stacks-queues:stack}{in-stack} and \omterm{def:g12:lists-stacks-queues:stack}{pushes}
it on the \omterm{def:g12:lists-stacks-queues:stack}{out-stack}, which reverses their order and puts the oldest on top.

\begin{omfigure}
\begin{tikzpicture}[font=\footnotesize, el/.style={draw, minimum width=8mm, minimum height=5mm, fill=omDef!10}]
  \foreach \x/\lab in {0/{in}, 2/{out}, 7/{in}, 9/{out}} {
    \draw[thick] (\x-0.5,1.9) -- (\x-0.5,-0.05) -- (\x+0.5,-0.05) -- (\x+0.5,1.9);
    \node[below] at (\x,-0.1) {\lab};
  }
  \foreach \v [count=\i] in {1, 2, 3} \node[el] at (0,{0.25+0.5*(\i-1)}) {$\v$};
  \foreach \v [count=\i] in {3, 2} \node[el] at (9,{0.25+0.5*(\i-1)}) {$\v$};
  \node[el, fill=omThm!15] (d) at (10.9,1.0) {$1$};
  \node[font=\scriptsize, omThm, above] at (d.north) {dequeued};
  \draw[->, omThm, thick] (9,1.2) to[bend left] (d.west);
  \draw[->, thick] (3.1,1) -- node[above, align=center, font=\scriptsize] {dequeue: out is\\ empty, so transfer}
    (5.9,1);
  \node[font=\scriptsize, black!60] at (1,2.3) {after enqueueing $1, 2, 3$};
  \node[font=\scriptsize, black!60] at (8,2.3) {after the first dequeue};
\end{tikzpicture}

{\small A \omterm{def:g12:lists-stacks-queues:queue}{queue} from two \omterm{def:g12:lists-stacks-queues:stack}{stacks}. The transfer \omterm{def:g12:lists-stacks-queues:stack}{pops} $3$, $2$, $1$ from the \omterm{def:g12:lists-stacks-queues:stack}{in-stack} and \omterm{def:g12:lists-stacks-queues:stack}{pushes} them on the \omterm{def:g12:lists-stacks-queues:stack}{out-stack},
where $1$, the oldest, is now on top. The next two \omterm{def:g12:lists-stacks-queues:queue}{dequeues} \omterm{def:g12:lists-stacks-queues:stack}{pop} $2$ and $3$ with no transfer.}
\end{omfigure}

\omcode{code/grade-12/03-lists-stacks-queues/listings/queue_two_stacks.py}{1}{26}{A \omterm{def:g12:lists-stacks-queues:queue}{queue} made of two
\omterm{def:g12:lists-stacks-queues:stack}{stacks}.}\label{lst:g12:lists-stacks-queues:queue-stacks}

One \omterm{def:g12:lists-stacks-queues:queue}{dequeue} can be slow: it may transfer every element. But an element is transferred at most once, since it never
goes back to the \omterm{def:g12:lists-stacks-queues:stack}{in-stack}. Counting a ``move'' each time an element is pushed on the \omterm{def:g12:lists-stacks-queues:stack}{in-stack}, transferred, or popped
from the \omterm{def:g12:lists-stacks-queues:stack}{out-stack}, each element makes at most three moves in its whole life: $n$ \omterm{def:g12:lists-stacks-queues:queue}{enqueues} followed by any \omterm{def:g12:lists-stacks-queues:queue}{dequeues}
cost at most $3n$ moves in all, however they are interleaved (exercise \ref{exo:g12:lists-stacks-queues:12}).

The two \omterm{def:g12:lists-stacks-queues:queue}{queues} pass one suite of tests, written like the \omterm{def:g12:lists-stacks-queues:stack}{stack} suite; here is one of its four tests.

\omcode{code/grade-12/03-lists-stacks-queues/listings/queues_tests.py}{18}{25}{A test that mixes \omterm{def:g12:lists-stacks-queues:queue}{enqueues} and
\omterm{def:g12:lists-stacks-queues:queue}{dequeues}.}\label{lst:g12:lists-stacks-queues:queues-tests}

\omoutput{code/grade-12/03-lists-stacks-queues/outputs/queues_tests.txt}{Both \omterm{def:g12:lists-stacks-queues:queue}{queues} pass.}%
\label{out:g12:lists-stacks-queues:queues-tests}

\section{Using them}\label{sec:g12:lists-stacks-queues:using}

\subsection*{Checking brackets}

In a program, every opening bracket must be closed by a bracket of the same kind, the last opened being the first
closed: a \omterm{def:g12:lists-stacks-queues:stack}{stack}. Each opening is pushed; each closing must match the top, which is popped; at the end nothing may be
left open.

\omcode{code/grade-12/03-lists-stacks-queues/listings/brackets.py}{1}{18}{Balanced brackets, with a
\omterm{def:g12:lists-stacks-queues:stack}{stack}.}\label{lst:g12:lists-stacks-queues:brackets}

\omoutput{code/grade-12/03-lists-stacks-queues/outputs/brackets.txt}{\code{(]} closes the wrong kind; \code{((x)}
leaves one open; \code{x)(} closes before opening.}\label{out:g12:lists-stacks-queues:brackets}

\subsection*{Computing in postfix notation}

\begin{definition}[Postfix notation]\label{def:g12:lists-stacks-queues:postfix}
In \emph{postfix notation}\index{postfix notation}, an operator is written after its two operands: $3 + 4$ is
\code{3 4 +}, and $(3 + 4) \times 2$ is \code{3 4 + 2 *}. No brackets and no priority rules are needed: the order of
the operators says everything.
\end{definition}

A postfix expression is evaluated with one \omterm{def:g12:lists-stacks-queues:stack}{stack}: a number is pushed; an operator \omterm{def:g12:lists-stacks-queues:stack}{pops} two numbers, the right operand
first, and \omterm{def:g12:lists-stacks-queues:stack}{pushes} the result. At the end the \omterm{def:g12:lists-stacks-queues:stack}{stack} holds exactly one number, the value.

\begin{omfigure}
\begin{tikzpicture}[font=\footnotesize, el/.style={draw, minimum width=8mm, minimum height=5mm, fill=omDef!10}]
  \foreach \t/\x in {3/0, 4/1.6, +/3.2, 2/4.8, */6.4} \node at (\x,-0.5) {read \texttt{\t}};
  \node[el] at (0,0.25) {$3$};
  \node[el] at (1.6,0.25) {$3$}; \node[el] at (1.6,0.75) {$4$};
  \node[el, fill=omThm!15] at (3.2,0.25) {$7$};
  \node[el] at (4.8,0.25) {$7$}; \node[el] at (4.8,0.75) {$2$};
  \node[el, fill=omThm!15] at (6.4,0.25) {$14$};
  \foreach \x in {0, 1.6, 3.2, 4.8, 6.4} \draw[thick] (\x-0.5,1.3) -- (\x-0.5,-0.05) -- (\x+0.5,-0.05) -- (\x+0.5,1.3);
\end{tikzpicture}

{\small The \omterm{def:g12:lists-stacks-queues:stack}{stack} while \code{3 4 + 2 *} is read, one picture per token. An operator replaces the two numbers on top
by its result (shaded); $14$ is left alone at the end.}
\end{omfigure}

\omcode{code/grade-12/03-lists-stacks-queues/listings/postfix.py}{1}{30}{A postfix calculator.}%
\label{lst:g12:lists-stacks-queues:postfix}

\omoutput{code/grade-12/03-lists-stacks-queues/outputs/postfix.txt}{The top after each token, then three values and a
malformed expression.}\label{out:g12:lists-stacks-queues:postfix}

\subsection*{A waiting line}

A counter serves people in the order they came: a \omterm{def:g12:lists-stacks-queues:queue}{queue}. Here people arrive at random, and the counter serves one
every two minutes.

\omcode{code/grade-12/03-lists-stacks-queues/listings/waiting_line.py}{1}{13}{Ten minutes at a
counter.}\label{lst:g12:lists-stacks-queues:line}

\omoutput{code/grade-12/03-lists-stacks-queues/outputs/waiting_line.txt}{Each person leaves in the order of arrival;
\code{c31} waits four minutes behind \code{c30}.}\label{out:g12:lists-stacks-queues:line}

\subsection*{Choosing a structure}

The choice follows from the operations the program needs. Undo, brackets, the calls of a \omterm{def:g12:recursion:recursion}{recursive function}: a
\omterm{def:g12:lists-stacks-queues:stack}{stack}. Services in order of arrival, a printer, the visit of a \omterm{def:g10:graphs-first-look:graph}{graph} level by level (chapter
\ref{ch:g12:graph-algorithms}): a \omterm{def:g12:lists-stacks-queues:queue}{queue}. Many insertions at the front: a \omterm{def:g12:lists-stacks-queues:linked-list}{linked list}. Reading by index: a Python list.
Finding a value from a key: a dictionary, which does not scan its keys at all.

\omcode{code/grade-12/03-lists-stacks-queues/listings/list_vs_dict.py}{1}{20}{Finding a key in a list of pairs and in
a dictionary.}\label{lst:g12:lists-stacks-queues:list-dict}

\omoutput{code/grade-12/03-lists-stacks-queues/outputs/list_vs_dict.txt}{On the author's laptop: the scan grows ten
times with $n$, the dictionary stays at about $0.01$ microseconds.}\label{out:g12:lists-stacks-queues:list-dict}

\section{Exercises}\label{sec:g12:lists-stacks-queues:exercises}

\begin{exercise}[$\star$]\label{exo:g12:lists-stacks-queues:1}
From an empty \omterm{def:g12:lists-stacks-queues:stack}{stack}: \omterm{def:g12:lists-stacks-queues:stack}{push} $1$, \omterm{def:g12:lists-stacks-queues:stack}{push} $2$, \omterm{def:g12:lists-stacks-queues:stack}{push} $3$, \omterm{def:g12:lists-stacks-queues:stack}{pop}, \omterm{def:g12:lists-stacks-queues:stack}{push} $4$, \omterm{def:g12:lists-stacks-queues:stack}{pop}, \omterm{def:g12:lists-stacks-queues:stack}{pop}. What did the \omterm{def:g12:lists-stacks-queues:stack}{pops} return, and what is
left?
\end{exercise}

\begin{exercise}[$\star$]\label{exo:g12:lists-stacks-queues:2}
From an empty \omterm{def:g12:lists-stacks-queues:queue}{queue}: \omterm{def:g12:lists-stacks-queues:queue}{enqueue} \code{a}, \code{b}, \code{c}, \omterm{def:g12:lists-stacks-queues:queue}{dequeue}, \omterm{def:g12:lists-stacks-queues:queue}{enqueue} \code{d}, \omterm{def:g12:lists-stacks-queues:queue}{dequeue}. What did the \omterm{def:g12:lists-stacks-queues:queue}{dequeues}
return, and what is left, front first?
\end{exercise}

\begin{exercise}[$\star$]\label{exo:g12:lists-stacks-queues:3}
Give the \omterm{def:g12:lists-stacks-queues:stack}{stack} after each character of \code{{[()]}(}, and the answer of \code{balanced}.
\end{exercise}

\begin{exercise}[$\star$]\label{exo:g12:lists-stacks-queues:4}
Evaluate \code{5 1 2 + 4 * + 3 -} by hand, with the \omterm{def:g12:lists-stacks-queues:stack}{stack} after each token.
\end{exercise}

\begin{exercise}[$\star$]\label{exo:g12:lists-stacks-queues:5}
Add a method \code{size} to both \omterm{def:g12:lists-stacks-queues:stack}{stacks}. In the linked \omterm{def:g12:lists-stacks-queues:stack}{stack}, it must not walk the \omterm{def:g12:lists-stacks-queues:linked-list}{cells}.
\end{exercise}

\begin{exercise}[$\star\star$]\label{exo:g12:lists-stacks-queues:6}
If \code{pop} did not test for an empty \omterm{def:g12:lists-stacks-queues:stack}{stack}, what would each \omterm{def:g11:modules-and-interfaces:interface}{implementation} raise on an empty \omterm{def:g12:lists-stacks-queues:stack}{stack}? Why is the
test worth writing?
\end{exercise}

\begin{exercise}[$\star\star$]\label{exo:g12:lists-stacks-queues:7}
Write the length of a \omterm{def:g12:lists-stacks-queues:linked-list}{linked list} given by its head, recursively and with a loop.
\end{exercise}

\begin{exercise}[$\star\star$]\label{exo:g12:lists-stacks-queues:8}
A student's \omterm{def:g12:lists-stacks-queues:queue}{queue} decides that it is empty by \code{self._tail is None}, and its \code{dequeue} never changes
\code{_tail}. \omterm{def:g12:lists-stacks-queues:queue}{Enqueue} $1$, \omterm{def:g12:lists-stacks-queues:queue}{dequeue}, \omterm{def:g12:lists-stacks-queues:queue}{enqueue} $2$: where is $2$, and what does the next \omterm{def:g12:lists-stacks-queues:queue}{dequeue} do?
\end{exercise}

\begin{exercise}[$\star\star$]\label{exo:g12:lists-stacks-queues:9}
On the \omterm{def:g12:lists-stacks-queues:stack}{two-stack} \omterm{def:g12:lists-stacks-queues:queue}{queue}: \omterm{def:g12:lists-stacks-queues:queue}{enqueue} $1$, $2$, $3$, \omterm{def:g12:lists-stacks-queues:queue}{dequeue}, \omterm{def:g12:lists-stacks-queues:queue}{enqueue} $4$, then \omterm{def:g12:lists-stacks-queues:queue}{dequeue} three times. Count the moves after
each operation.
\end{exercise}

\begin{exercise}[$\star\star$]\label{exo:g12:lists-stacks-queues:10}
Write \code{reverse(head)}, which reverses a \omterm{def:g12:lists-stacks-queues:linked-list}{linked list} by changing its links, without creating a cell, and returns
the new head.
\end{exercise}

\begin{exercise}[$\star\star$]\label{exo:g12:lists-stacks-queues:11}
Write $(2 + 3) \times (7 - 4) / 5$ in \omterm{def:g12:lists-stacks-queues:postfix}{postfix notation}, and evaluate it with the \omterm{def:g12:lists-stacks-queues:stack}{stack} method.
\end{exercise}

\begin{exercise}[$\star\star\star$]\label{exo:g12:lists-stacks-queues:12}
Prove that, in the \omterm{def:g12:lists-stacks-queues:stack}{two-stack} \omterm{def:g12:lists-stacks-queues:queue}{queue}, each element makes at most three moves, and deduce that any sequence of
operations containing $n$ \omterm{def:g12:lists-stacks-queues:queue}{enqueues} makes at most $3n$ moves.
\end{exercise}

\begin{exercise}[$\star\star\star$]\label{exo:g12:lists-stacks-queues:13}
Write a \omterm{def:g12:lists-stacks-queues:queue}{queue} of fixed capacity kept in a Python list of that length, with the index of the oldest element and the
number of elements; the indices wrap around to $0$ at the end of the list. A full \omterm{def:g12:lists-stacks-queues:queue}{queue} refuses to \omterm{def:g12:lists-stacks-queues:queue}{enqueue}.
\end{exercise}

\begin{exercise}[$\star\star\star$]\label{exo:g12:lists-stacks-queues:14}
Write \code{to_postfix(expression)} for expressions of numbers, \code{+}, \code{-}, \code{*}, \code{/} and brackets,
tokens separated by spaces. Numbers go straight to the output; operators wait on a \omterm{def:g12:lists-stacks-queues:stack}{stack}, which first gives up the
operators of priority greater than or equal to the new one; a closing bracket empties the \omterm{def:g12:lists-stacks-queues:stack}{stack} down to its opening
bracket.
\end{exercise}

\begin{exercise}[$\star\star\star$]\label{exo:g12:lists-stacks-queues:15}
Call a string of the characters \code{(}, \code{)}, \code{[}, \code{]} \emph{balanced} if it is empty, or of the form
$(u)v$ or $[u]v$ with $u$ and $v$ balanced. Prove that \code{balanced} returns \code{True} exactly on the balanced
strings.
\end{exercise}

\section{Problem: The Queue at the Post Office}\label{sec:g12:lists-stacks-queues:problem}

\begin{problem}[{Weekend problem --- a day at the counter}]\label{pb:g12:lists-stacks-queues:1}
A post office opens for $480$ minutes. In each minute, a customer arrives with chance $0.3$; one clerk serves the
line in order, each service taking $1$ to $5$ minutes. At closing time the doors shut and the people inside are still
served. The program runs the same day, with the same random numbers, on both \omterm{def:g12:lists-stacks-queues:queue}{queues}.

\omcode{code/grade-12/03-lists-stacks-queues/listings/post_office.py}{1}{33}{A day at the post office, on both
\omterm{def:g12:lists-stacks-queues:queue}{queues}.}\label{lst:g12:lists-stacks-queues:post}

\omoutput{code/grade-12/03-lists-stacks-queues/outputs/post_office.txt}{One day, two \omterm{def:g11:modules-and-interfaces:interface}{implementations}.}%
\label{out:g12:lists-stacks-queues:post}

\needspace{4\baselineskip}\medskip\noindent\textbf{Part I --- The queue as an abstract data type.}
\begin{enumerate}
  \item Which operations of the \omterm{def:g12:lists-stacks-queues:queue}{queue} does \code{simulate} use? Which attribute does it read that is not part of
    the \omterm{def:g12:lists-stacks-queues:adt}{abstract data type}?
  \item Why can \code{simulate} take the class \code{Queue} as an argument and work with either?
  \item What does the \omterm{def:g12:lists-stacks-queues:queue}{queue} hold: the customers' names, or something else? Why is that enough?
  \item The variable \code{waiting} counts the people in line. Why does the program keep it, rather than ask the
    \omterm{def:g12:lists-stacks-queues:queue}{queue}?
\end{enumerate}

\needspace{4\baselineskip}\medskip\noindent\textbf{Part II --- Two implementations.}
\begin{enumerate}[resume]
  \item In \code{queue_linked}, which links change at an \omterm{def:g12:lists-stacks-queues:queue}{enqueue} on an empty \omterm{def:g12:lists-stacks-queues:queue}{queue}, and on a \omterm{def:g12:lists-stacks-queues:queue}{queue} of three?
  \item Why must \code{dequeue} set \code{_tail} to \code{None} when the \omterm{def:g12:lists-stacks-queues:queue}{queue} becomes empty?
  \item In \code{queue_two_stacks}, when does a \omterm{def:g12:lists-stacks-queues:queue}{dequeue} transfer elements, and how many?
  \item Give a sequence of $6$ operations in which one \omterm{def:g12:lists-stacks-queues:queue}{dequeue} makes $4$ moves.
  \item Which \omterm{def:g11:modules-and-interfaces:interface}{implementation} would you pick for a \omterm{def:g12:lists-stacks-queues:queue}{queue} that is mostly empty? For one where every \omterm{def:g12:lists-stacks-queues:queue}{dequeue} must be
    fast?
\end{enumerate}

\needspace{4\baselineskip}\medskip\noindent\textbf{Part III --- The simulation.}
\begin{enumerate}[resume]
  \item Why is the random generator built from a fixed seed?
  \item Why do both \omterm{def:g11:modules-and-interfaces:interface}{implementations} give exactly the same customers, line and waits?
  \item A service lasts $3$ minutes on average over the five durations, and a customer arrives about every
    $1/0.3 \approx 3.3$ minutes. What happens to the line if the services become a little longer?
  \item Why does the loop go on after minute $480$, and when does it stop?
  \item Explain how the program finds each customer's wait.
\end{enumerate}

\needspace{4\baselineskip}\medskip\noindent\textbf{Part IV --- Results and costs.}
\begin{enumerate}[resume]
  \item Read the number of customers, the longest line and the longest wait.
  \item How many moves does the linked \omterm{def:g12:lists-stacks-queues:queue}{queue} make per customer? Why exactly that number?
  \item How many does the \omterm{def:g12:lists-stacks-queues:stack}{two-stack} \omterm{def:g12:lists-stacks-queues:queue}{queue} make per customer? Why exactly $3$, and not at most $3$?
  \item Is the \omterm{def:g12:lists-stacks-queues:stack}{two-stack} \omterm{def:g12:lists-stacks-queues:queue}{queue} slower in time? Compare the work of one move in each.
  \item What would change in the moves if the program stopped at minute $480$, with people still in line?
  \item State the result of the day, and of the comparison.
\end{enumerate}
\end{problem}
